\section{Existence requires a witness and evidence}
To prove $\exists n\in\NN:n+1=4$, supply the candidate $3$ and evidence that it satisfies the property:
\par\noindent\begin{minipage}{\linewidth}
\begin{lstlisting}
theorem a_witness : Exists (fun n : Nat => n + 1 = 4) := by
  exact Exists.intro 3 rfl
\end{lstlisting}
\end{minipage}\par

The expression \code{Exists (fun n : Nat => n + 1 = 4)} states that some natural number has the displayed property. This is the typed form of the existential quantifier. \code{Exists.intro} packages a witness with its proof. The equality is computational in this case, so \code{rfl} suffices. For a more difficult witness, the second argument would be a substantive proof.

Conversely, using an existential hypothesis requires unpacking both components. Knowing merely that ``something exists'' is not enough to apply a formula to it without introducing a name and its guaranteed property. Chapter~16 uses this operation in a complete surjectivity proof.

\subsection*{The type of a number changes available operations}
A natural number has no negative values. Lean's natural-number subtraction is truncated at zero: \code{(7 - 10 : Nat)} is zero. Integer subtraction is different: \code{(7 - 10 : Int)} is minus three. The type annotation inside the parentheses specifies which operation is meant.

In ordinary mathematics we also need this distinction, although informal notation sometimes leaves the domain implicit. An argument that subtracts a larger integer from a smaller one and then treats the result as a negative number cannot be translated unchanged into natural-number subtraction. Before asking an assistant to repair a formal proof, check whether it has silently changed the kind of number being used.
